\begin{center}\textls[120]{c) Konforme Abbildung. Geometrisches.}\end{center}
\noindent\hspace*{\parindent}In dem oben (S. 167) erwähnten Briefe vom 8. Juni 1816 fragt \textsc{Schumacher}
bei \textsc{Gauss} an\footnote{Vergl. \textsc{Galle}, Werke XI, 2, Abh. 1, S. 32 ff., wo auch die betreffende Briefstelle wiedergegeben ist.}, ob es wohl möglich wäre, die für Dänemark ihm
aufgetragene Gradmessung\footnote{Vergl. über den auf \textsc{Schumacher}s Initiative geschaffenen geodätischen Dienst in Dänemark \textsc{N. E. Nörlund}, \textit{Notice historique etc.} in den Verhandlungen der 5. Tagung der baltischen geodätischen Kommission, Helsingfors, 1931, S. 6 ff.} durch Hannover fortzusetzen. \textsc{Gauss} antwortet
am 5. Juli\footnote{Siehe die Briefstelle bei \textsc{Galle}, a. a. O. S. 33 ff.}, er bezweifle nicht, dass dies, wenn auch nicht sogleich, so
doch in Zukunft möglich sein werde, und dass er eine eigene Methode besitze,
»um die gemessenen Dreiecke im Kalkul zu behandeln«. Aus dieser
Korrespondenz nimmt \textsc{Gauss}' Tätigkeit auf dem Gebiete der Geodäsie ihren
Ursprung, eine Tätigkeit, die sich mit wechselnder Intensität über einen Zeitraum
\par\clearpage\SourceLeaf{256}\ArithmeticHead{171}{R}{VIIIc. KONFORME ABBILDUNG. GEOMETRISCHES.}
\noindent von mehr als dreissig Jahren erstreckt. Charakteristisch für die Art,
wie sich bei \textsc{Gauss} aus praktischen Anregungen mathematische Ideen von der
grössten Tragweite entwickeln, ist die folgende Stelle in demselben Briefe
von \textsc{Gauss} an \textsc{Schumacher} (Werke VIII, S. 370): »Mit \textsc{Lindenau} habe ich
auch über eine Preisfrage konferiert, die in der neuen Zeitschrift mit dem
Preise von 100 Dukaten aufgegeben werden soll. Mir war eine interessante
Aufgabe eingefallen, nämlich:
\begin{quote}
»Allgemein eine gegebene Fläche so auf einer andern (gegebenen) zu
projizieren (abzubilden), dass das Bild dem Originale in den kleinsten
Teilen ähnlich werde«\footnote{Näheres über die Korrespondenz mit \textsc{Lindenau} bei \textsc{Galle} a. a. O. S. 41 und 50—53.}.
\end{quote}
\noindent Als Preisfrage der Kopenhagener Sozietät wurde diese Aufgabe 1820 veröffentlicht\footnote{Siehe Brief \textsc{Schumacher}s vom 20. Januar 1821, \textit{Briefwechsel Gauss-Schumacher} I, S. 202.}.
Nachdem sie 1822 wiederholt worden war, wendet sich \textsc{Gauss}
ihrer Bearbeitung zu\footnote{\textsc{Gauss} an \textsc{Schumacher}, 25. November 1822, ebenda S. 293.}. Die Abhandlung mit dem Motto »Ab his via sternitur
ad maiora« erscheint 1825\footnote{\textit{Allgemeine Auflösung der Aufgabe, die Teile einer gegebenen Fläche auf einer andern gegebenen Fläche so abzubilden, dass die Abbildung dem Abgebildeten in den kleinsten Teilen ähnlich wird}, Astronomische Abhandlungen 3, Werke IV, S. 189.}, sie ist vom Dezember 1822 datiert und
in deutscher Sprache verfasst\footnote{Vergl. die Briefstellen \textsc{Gauss-Schumacher}, a. a. O. S. 296, 297, 299, 300, 303, 317, 328, 335, ferner Werke XII, S. 222, 223 und \textsc{Stäckel}, Abh. 4 dieses Bandes, S. 84 ff.}. Die Wurzeln des Gedankenganges, der
\textsc{Gauss} zu dieser Aufgabe geführt hat, lassen sich weit zurück verfolgen\footnote{Vergl. \textsc{Galle} a. a. O. S. 53 und \textsc{Stäckel} a. a. O.}.
Dass jede funktionale Beziehung zwischen zwei komplexen Veränderlichen als
eine Abbildung (\textsc{Gauss} sagt oft Darstellung) der beiden Ebenen aufeinander
aufgefasst werden kann, war, wie wir in den vorhergehenden Abschnitten gesehen
haben, \textsc{Gauss} von jeher geläufig. Auch die besondere Eigenschaft einer
solchen Abbildung, in den kleinsten Teilen ähnlich zu sein, wird er wohl
frühzeitig erkannt haben. Im Art. 8 der Preisarbeit (Werke IV, S. 200) setzt
er, wenn die Beziehung \(X+Yi=f(x+yi)\) gilt, \(df(z)/dz=\varphi(z)\) und
\[
\varphi(x+yi)=\xi+\eta i=\sigma e^{\gamma i},\quad dX+idY=dS\mathbin{.}e^{Gi},\quad dx+idy=ds\mathbin{.}e^{gi},
\]
so dass also \(dS,ds\) die Linienelemente der beiden Ebenen sind, und folgert
\[
dX=\xi dx-\eta dy,\quad dY=\eta dx+\xi dy,
\]
\HistoricalFoot{}{22\textsuperscript{*}}
\par\clearpage\SourceLeaf{257}\ArithmeticHead{172}{L}{L. SCHLESINGER, ÜBER GAUSS' FUNKTIONENTHEORETISCHE ARBEITEN}
\noindent also die sogenannten \textsc{Cauchy-Riemann}schen partiellen Differentialgleichungen,
ferner \(dS=\sigma ds\), d. h. die Ähnlichkeit in den kleinsten Teilen, und \(G=g+\gamma\),
d. h. die Erhaltung der Winkel, u. z. mit Bewahrung des Drehungssinns. Nun
enthält schon der Art. 10 der \textit{Theoria attractionis} von 1813 (Werke V, S. 14)
die Bemerkung, dass es für die Darstellung einer Fläche besser sei, statt einer
Gleichung \(W(x,y,z)=0\) zwischen den rechtwinkligen Koordinaten »zwei
neue Veränderliche \(p,q\) einzuführen, derart, dass \(x,y,z\) als Funktionen dieser
Variabeln zu betrachten sind«. Interpretiert man dann \(p,q\) als rechtwinklige
Koordinaten in einer Ebene, sind ferner \(x,y,z\) eindeutig in \(p,q\) und
umgekehrt auch \(p,q\) durch \(x,y,z\) eindeutig bestimmt, so entspricht jedem
Punkte der Fläche ein Punkt der Ebene oder eines Teiles der Ebene und
umgekehrt, so dass die Ebene oder jener Teil derselben »gewissermassen ein
Bild der Fläche darstellt«. Damit ist natürlich auch die Abbildung von zwei
beliebigen Flächen aufeinander gegeben. Sind für die beiden Flächen die
rechtwinkligen Koordinaten \(x,y,z\) bezw. \(X,Y,Z\) Funktionen der beiden Parameter
\(t,u\) bezw. \(T,U\), so wird nach Art. 3 der Preisarbeit eine Abbildung
der ersten Fläche auf die zweite vollzogen sein, wenn man »\(T,U\) bestimmten
Funktionen von \(t,u\) gleichsetzt. Insofern die Abbildung gewissen Bedingungen
Genüge leisten soll, werden diese Funktionen nicht mehr willkürlich sein
dürfen«. In einer Notiz auf Seite 70 des Handbuchs 16, Bb (Werke VIII,
S. 37), die wohl gleichzeitig mit dem Briefe an \textsc{Schumacher}, also um die
Mitte des Jahres 1816 geschrieben sein dürfte, gibt \textsc{Gauss} »die allgemeine
Lösung der Aufgabe, eine krumme Fläche in den kleinsten Teilen ähnlich
auf eine Ebene zu entwerfen«. In den Art. 4 und 5 der Preisarbeit, wird
die Aufgabe der konformen\footnote{Die Bezeichnung \textls[120]{konform} für in den kleinsten Teilen ähnlich führt \textsc{Gauss} erst in den \textit{Untersuchungen über Gegenstände der höheren Geodäsie} I, 1843/44 ein; siehe Werke IV, S. 262 und 348.} Abbildung zweier Flächen aufeinander wie folgt
gelöst: Die Quadrate der Linienelemente der beiden Flächen sind
\[
\omega=e\,dt^2+2f\,dt\,du+g\,du^2,\quad \Omega=E\,dt^2+2F\,dt\,du+G\,du^2,
\]
wo \(e,f,g\) bezw. \(E,F,G\) in bekannter Weise aus den partiellen Derivierten
der Koordinaten nach den Parametern \(t,u\) zusammengesetzt sind\footnote{In der Preisarbeit kommen die Bezeichnungen \(E,F,G\) noch nicht vor, \textsc{Gauss} schreibt dort immer die vollständigen Ausdrücke hin; wir gebrauchen der Kürze halber die später (1827) von \textsc{Gauss} benutzten Zeichen.}. Die
\par\clearpage\SourceLeaf{258}\ArithmeticHead{173}{R}{VIIIc. KONFORME ABBILDUNG. GEOMETRISCHES.}
\noindent Ähnlichkeit in den kleinsten Teilen erfordert:
\begingroup\makeatletter\tagsleft@true\let\veqno\leqno\makeatother
\begin{equation*}
\frac Ee=\frac Ff=\frac Gg=m^2,\tag{10}
\end{equation*}
\endgroup
\noindent wo \(m\) »eine endliche Funktion von \(t,u\) sein wird; es drückt dann \(m\) das Verhältnis
aus, in welchem die Lineargrössen auf der ersten Fläche in ihrem
Abbild auf der zweiten vergrössert oder verkleinert werden. \mbox{.\,.\,.} In dem speziellen
Falle, wo \(m=1\) ist, wird \mbox{.\,.\,.} die eine Fläche sich auf der andern
abwickeln lassen«.

\textsc{Gauss} betrachtet dann im Art. 5 die Gleichung \(\omega=0\) der jetzt sogenannten
Minimalkurven auf der ersten Fläche, d. h.
\[
e\,dt+\left(f\pm i(eg-f^2)^{\frac12}\right)du=0,\quad eg-f^2>0,
\]
die integriert Lösungen von der Form \(p\pm qi=\mathrm{const}\) gibt, so dass
\[
\omega=n(dp^2+dq^2)
\]
wird, wo \(p,q,n\) reelle Funktionen von \(t,u\) sind. Hat man analog für die
zweite Fläche \(\Omega=N(dP^2+dQ^2)\), so erfordert die Bedingung (10), dass entweder
\[
\begin{array}{lrl}
&1)&P+Qi=\mathrm{const},\ \text{wenn }p+qi=\mathrm{const},\\
\text{oder}\quad&2)&P+Qi=\mathrm{const},\ \text{wenn }p-qi=\mathrm{const},
\end{array}
\]
d. h. aber, dass im ersten Falle \(P+Qi\) »bloss Funktion von \(p+qi\) und
ebenso \(P-Qi\) Funktion von \(p-qi\)«, im andern Falle \(P+Qi\) Funktion
von \(p-qi\) und \(P-Qi\) Funktion von \(p+qi\) sein wird\footnote{Zu dieser Definition der Funktion vergl. \textsc{Euler}, \textit{Calculus Integralis} III, Calc. Var. § 96, Opera omnia ser. I, vol. XIII, S. 412; über die konforme Abbildung bei \textsc{Gauss}' Vorgängern \textsc{Lambert}, \textsc{Euler} und \textsc{Lagrange} vergl. man \textsc{Stäckel}, Abh. 4 dieses Bandes, S. 92 ff.}. Die beiden Fälle
unterscheiden sich dadurch, dass (Art. 14, Werke IV, S 211) »stets bei dem
einen die Teile der Darstellung ähnliche Lage haben wie im Dargestellten,
bei dem andern \mbox{.\,.\,.} dagegen verkehrt liegen.« Diese Unterscheidung bedingt
aber, dass »an jeder der beiden Flächen zwei Seiten unterschieden werden,
wovon die eine als die obere, die andere als die untere betrachtet wird. Da
dieses etwas willkürliches ist, so sind beide Auflösungen gar nicht wesentlich
verschieden. \mbox{.\,.\,.} Bei unserer Auflösung konnte diese Unterscheidung gar nicht
vorkommen, da die Flächen bloss durch die Koordinaten ihrer Punkte bestimmt
waren«. Für den Fall zweier Ebenen (Art. 8) \(z=0,\ Z=0\) gibt die
\par\clearpage\SourceLeaf{259}\ArithmeticHead{174}{L}{L. SCHLESINGER, ÜBER GAUSS' FUNKTIONENTHEORETISCHE ARBEITEN.}
\noindent Auflösung \(X+Yi=f(x+yi)\) die ähnliche Lage (mit Erhaltung des Drehungssinns),
wenn man die beiden Koordinatensysteme \(x,y,z\) und \(X,Y,Z\) ähnlich
liegend nimmt, und bei beiden Ebenen die Seite, nach der hin \(z\) bezw. \(Z\)
positiv werden, als die obere betrachtet, während dann bei der Auflösung
\(X+Yi=f(x-yi)\) der Sinn sich umkehrt, so »dass dort rechts liegt, was
hier links ist«. Für zwei beliebige, durch die Gleichungen \(\psi(x,y,z)=0\),
\(\Psi(x,y,z)=0\) gegebene Flächen können die oberen Seiten auch etwa durch
\(\psi>0,\ \Psi>0\) charakterisiert werden, und \textsc{Gauss} zeigt dann, dass mit Beibehaltung
der oben benutzten Bezeichnungen, für \(P+Qi=f(p+qi)\) ähnliche
Lage stattfindet, wenn von den vier Determinanten
\[
\begin{gathered}
\Sigma\pm\frac{\partial\psi}{\partial x}\frac{\partial y}{\partial t}\frac{\partial z}{\partial u},\quad
\frac{\partial p}{\partial t}\frac{\partial q}{\partial u}-\frac{\partial p}{\partial u}\frac{\partial q}{\partial t},\\
\Sigma\pm\frac{\partial\Psi}{\partial X}\frac{\partial Y}{\partial T}\frac{\partial Z}{\partial U},\quad
\frac{\partial P}{\partial T}\frac{\partial Q}{\partial U}-\frac{\partial P}{\partial U}\frac{\partial Q}{\partial T}
\end{gathered}
\]
keine oder eine gerade Anzahl negativ ist, dagegen verkehrte Lage, wenn
dies für eine ungerade Anzahl eintritt. Nennt man, mit Beiseitelassung dieser
Unterscheidung, eine Funktion von \(p+qi\) bezw. \(P+Qi\) eine monogene Funktion
des Orts auf der ersten bezw. zweiten Fläche, so kann man (vergl. Art. 6)
sagen, dass durch die konforme Abbildung der beiden Flächen aufeinander,
die Gesamtheit der monogenen Funktionen des Orts auf der einen Fläche in
die entsprechende Gesamtheit auf der andern übergeführt wird. — Dabei ist
noch zu bemerken, dass den Erörterungen von \textsc{Gauss} die stillschweigende Annahme
zugrunde liegt, dass es sich stets um reguläre Stücke analytischer
Flächen handelt.

Zusammenfassend kann man sagen, dass \textsc{Gauss} die drei Typen der Abbildung
zweier Flächen aufeinander scharf erfasst und charakterisiert hat. Die
allgemeinste ist dadurch definiert, dass man »\(T,U\) bestimmten Funktionen von
\(t,u\) gleichsetzt«, wobei\begingroup\renewcommand{\thefootnote}{260a)}\footnote{Siehe den Brief an \textsc{Hansen} vom 11. Dezember 1825, Werke XII, S. 6.}\addtocounter{footnote}{-1}\endgroup\ »jedem Punkte der einen Fläche nach irgend einem
stetigen Gesetz ein Punkt der andern korrespondieren soll«; sie ist also das,
was wir heute als \textls[120]{topologische Abbildung}\linebreak bezeichnen. Es ist kein Anhaltspunkt
dafür vorhanden, dass \textsc{Gauss} etwas über die Invarianten einer
solchen Abbildung gewusst hat, obwohl er (vergl. \textsc{Stäckel}, Abhandl. 4 dieses
Bandes, S. 49 ff.) viel über die »Geometria situs« nachgedacht und, wie wir
sehen werden, auch die Bedeutung der Zusammenhangsverhältnisse schlichter
ebener Bereiche für die Integralrechnung im komplexen Gebiet erkannt hatte.
\par\clearpage\SourceLeaf{260}\ArithmeticHead{175}{R}{VIIIc. KONFORME ABBILDUNG. GEOMETRISCHES.}
\noindent Für den zweiten Typus der Abbildung, die \textls[120]{konforme}, hat er die allgemeine
Theorie aufgestellt und in der Gesamtheit der monogenen Funktionen des Orts
auf der Fläche die Invariante dieser Abbildung erkannt. Wie weit er in der
Behandlung spezieller Probleme der konformen Abbildung ebener Bereiche aufeinander
vorgedrungen ist, werden wir im folgenden Abschnitt sehen, hier können
wir vorerst auf die Beispiele der Preisarbeit (Art. 9—13) verweisen, von deren
erneuter Durcharbeitung vielleicht noch Anregungen zu erwarten sind. Endlich
— und das dürfte wohl \textsc{Gauss}' bedeutendste geometrische Erkenntnis gewesen
sein — hat er die Abwicklung oder, wie wir heute lieber sagen, die \textls[120]{isometrische}
Abbildung, nicht nur durch die Invarianz des Linienelements charakterisiert,
und damit die Unterscheidung zwischen inneren (absoluten) und von der
Einlagerung in den umgebenden Raum abhängenden Eigenschaften einer Fläche
gegeben, sondern er hat auch schon um 1816 (Werke VIII, S. 372, 385 aus
Handbuch 16, Bb, S. 71 bezw. 83) in der Gesamtkrümmung eines Flächenstücks,
in dem Krümmungsmass der Fläche in einem regulären Punkte und in den geodätischen
Linien, wozu etwas später (wohl um 1822) sich noch die Seitenkrümmung
(geodätische Krümmung) einer Kurve auf der Fläche (Werke VIII, S. 386)
hinzugesellte, solche innere oder absolute Eigenschaften aufgestellt. Jedenfalls
war um 1822 der wesentliche Inhalt der 1827 veröffentlichten \textit{Disquisitiones
generales circa superficies curvas} (Werke IV, S. 217) nicht nur entdeckt, sondern
auch zum Teil aufgezeichnet (Werke VIII, S. 374, datiert 13. Dezember 1822),
während aus dem Jahre 1825 eine ausführliche Redaktion in deutscher Sprache
(Werke VIII, S. 408—443) vorliegt.

Aber noch eine Reihe weiterer geometrischer Einsichten hat sich \textsc{Gauss}
um die Zeit von 1816 eröffnet. Aus einem Briefe von \textsc{F. L. Wachter} an
\textsc{Gauss} vom 12. Dezember 1816\footnote{Veröffentlicht durch \textsc{P. Stäckel}, Mathem. Annalen 54, 1901, S. 61—69, Teile daraus, Werke VIII, S. 175 und X, 1, S. 481, 482; vergl. zu dem im Text folgenden, \textsc{Stäckel}, Abh. 4 dieses Bandes, S. 28—31 und die Bemerkung Werke, X, 1, S. 481.} geht hervor, dass \textsc{Gauss} im Jahre 1816
einerseits im Besitze der nichteuklidischen (oder, wie er es nannte, \textls[120]{transzendenten})
Trigonometrie gewesen ist\footnote{Vergl. die aus späterer Zeit stammende Darstellung dieser Trigonometrie Werke VIII, S. 255 ff. und die zugehörige Bemerkung von \textsc{Stäckel}, ebenda S. 257—265, ferner Werke X, 1, S. 451—456.} und andererseits mit der Vorstellung
eines »Raumes von beliebig vielen Dimensionen« vertraut war\footnote{Vergl. hierzu die Auszüge aus einer Vorlesung von 1850/51 nach einer Ausarbeitung von \textsc{A. Ritter}, Werke X, 1, S. 472 und ebenda, S. 467 einen Auszug aus \textsc{Ritter}s Inauguraldissertation, Göttingen 1853; siehe auch \textsc{Stäckel}, Abh. 4 dieses Bandes, S. 117, 118.}. Auch geht
\par\clearpage\SourceLeaf{261}\ArithmeticHead{176}{L}{L. SCHLESINGER, ÜBER GAUSS' FUNKTIONENTHEORETISCHE ARBEITEN.}
\noindent aus einer \textit{Anzeige} (Werke IV, S. 364) vom 20. April 1816, sowie aus Briefen
an \textsc{Olbers} (vom 28. April 1817, Werke VIII, S. 177) und an \textsc{Gerling} (vom
11. April 1816, Werke VIII, S. 168 und vom 16. März 1819, ebenda S. 182)
hervor, dass er damals zur Überzeugung von der Unbeweisbarkeit des Parallelenaxioms
und der Widerspruchsfreiheit der nichteuklidischen Geometrie\footnote{Darunter ist hier nur die \textsc{Bolyai-Lobatschefskij}sche, sog. hyperbolische Geometrie zu verstehen; ob \textsc{Gauss} später auch die elliptisch-sphärische Geometrie gekannt hat, ist unsicher, ebenso wie die Entscheidung der Frage, ob er zu der Erkenntnis des Zusammenhangs zwischen der nichteuklidischen Geometrie der Ebene und der Geometrie auf einer Fläche von konstantem negativen Krümmungsmass vorgedrungen war; vergl. hierzu \textsc{Stäckel}, Abh. 4 dieses Bandes, S. 39, 40 und 108, 109.}
gelangt war.

Dass sich bei einem Denker wie \textsc{Gauss} diese geometrischen Einsichten
zu einer festbegründeten Auffassung von der Wesenheit des Raumes verdichtet
haben, ist nicht zu verwundern. In dem Briefe an \textsc{Hansen} vom 11. Dezember
1825 (Werke XII, S. 8) schreibt er: »Ich habe mich in diesem Herbst
sehr viel mit der allgemeinen Betrachtung der krummen Flächen beschäftigt,
welches in ein unabsehbares Feld führt. \mbox{.\,.\,.} Jene Untersuchungen greifen tief in
vieles andere, ich möchte sogar sagen, in die Metaphysik der Raumlehre ein,
und nur mit Mühe kann ich mich von solchen daraus entspringenden Folgen,
wie z. B. die wahre Metaphysik der negativen und imaginären Grössen ist,
losreissen. Der wahre Sinn des \(\sqrt{-1}\) steht mir dabei mit grosser Lebendigkeit
vor der Seele, aber es wird sehr schwer sein, ihn in Worte zu fassen\mbox{.\,.\,.}«
Das »Eingreifen in die Metaphysik der Raumlehre« wird man dahin zu deuten
haben, dass \textsc{Gauss} erkannte, dass der Raum ein besonderer Fall einer Mannigfaltigkeit
von drei Dimensionen ist, ebenso wie die Ebene ein besonderer Fall
der Fläche. Natürlich war ihm klar, dass die Punkte einer Mannigfaltigkeit
von drei Dimensionen ebenso durch drei Bestimmungsstücke (Koordinaten)
\(t,u,v\) gegeben werden, wie die Punkte einer Fläche durch zwei Grössen \(t,u\)
und die Punkte einer Linie durch eine \(t\). Von hier aus gesehen würde dann
die »wahre Metaphysik der negativen und imaginären Grössen« darin bestehen,
dass man die Punkte einer Linie den Werten einer reellen Veränderlichen \(t\),
also den positiven und negativen Grössen, die Punkte einer Fläche den komplexen
Grössen \(t+ui\) zuordnen kann, und dass für diese ein- und zweidimensionalen
Vektoren die gewöhnlichen Rechnungsregeln (insbesondere
\par\clearpage\SourceLeaf{262}\ArithmeticHead{177}{R}{VIIIc. KONFORME ABBILDUNG. GEOMETRISCHES.}
\noindent kommutative Multiplikation) gelten, während (wie es in der \textit{Anzeige} von 1831, Werke
II, S. 178, heisst) »die Relationen zwischen Dingen, die eine Mannigfaltigkeit
von mehr als zwei Dimensionen darbieten, nicht noch andere, in der allgemeinen
Arithmetik zulässige Grössen liefern können«. Und der »wahre Sinn
des \(\sqrt{-1}\)« dürfte darin zu erblicken sein, dass \textsc{Gauss} irgendwie zu der Einsicht
gelangt war, dass die zweidimensionalen Vektoren auch insofern eine
singuläre Stellung einnehmen, als einmal die Möglichkeit der konformen Abbildung
und andererseits auch die der Verbiegung oder isometrischen Abbildung
im allgemeinen auf Mannigfaltigkeiten von mehr als zwei Dimensionen
nicht übertragen werden kann\footnote{Siehe die neuere Literatur zu diesen beiden Fragen bei \textsc{Berwald}, Encyklopädie, Band III, 3, S. 158 und 167.}.

Wenn aus der eben besprochenen Briefstelle an \textsc{Hansen} die Auffassung hervorleuchtet,
dass der Raum eine dreifach ausgedehnte Zahlenmannigfaltigkeit sei,
so können wir auch Äusserungen von \textsc{Gauss} — allerdings aus etwas späterer
Zeit — anführen, die sich auf die Natur dieser Mannigfaltigkeit beziehen. In
zwei Briefen an \textsc{Bessel} vom 27. Januar 1829 und 9. April 1830 (Werke VIII,
S. 200, 201) heisst es: »Meine Überzeugung, dass wir die Geometrie nicht
vollständig a priori begründen können, ist womöglich noch fester geworden«
und »Nach meiner innigsten Überzeugung hat die Raumlehre in unserem
Wissen eine ganz andere Stellung als die reine Grössenlehre \mbox{.\,.\,.}; wir müssen
in Demut zugeben, dass, wenn die Zahl bloss unseres Geistes Produkt ist, der
Raum auch ausser unserem Geiste eine Realität hat, der wir a priori ihre
Gesetze nicht vollständig vorschreiben können.« Für diesen \textls[120]{empirischen}
Charakter des Raumes bringt \textsc{Gauss} im wesentlichen zwei Gründe vor.

Der erste findet sich in der bereits genannten \textit{Anzeige} von 1831 (Werke II,
S. 177), wo bei der geometrischen Repräsentation der komplexen Grössen gesagt
wird, man könne »den sich auf \(+i\) beziehenden Punkt nach Gefallen rechts
oder links« von der \(+1\) mit \(-1\) verbindenden Geraden nehmen, und es dann
in bezug auf den Unterschied zwischen rechts und links weiter heisst\begingroup\renewcommand{\thefootnote}{265a)}\footnote{Vergl. zum folgenden auch die aus 1850 stammende Fussnote, Werke X, 1, S. 408, 409, in der auch auf die \textit{Anzeige} von 1831 hingewiesen wird.}\addtocounter{footnote}{-1}\endgroup:
\par\noindent\hangindent=7.5em\hangafter=1
(1. \textls[120]{Bemerkung}) »Dieser Unterschied zwischen rechts und links ist, sobald
man vorwärts und rückwärts in der Ebene und oben und
\HistoricalFoot{X 2 Abh. 2.}{23}
\par\clearpage\SourceLeaf{263}\ArithmeticHead{178}{L}{L. SCHLESINGER, ÜBER GAUSS' FUNKTIONENTHEORETISCHE ARBEITEN.}
\noindent\hspace*{7.5em}\begin{minipage}[t]{\dimexpr\linewidth-7.5em\relax}
unten in Beziehung auf die beiden Seiten der Ebene einmal
nach Gefallen festgesetzt hat, in sich völlig bestimmt,
\end{minipage}
\par\noindent\hangindent=7.5em\hangafter=1
(2. \textls[120]{Bemerkung}) wenn wir gleich unsere Anschauung dieses Unterschiedes
andern nur durch Nachweisung an wirklich vorhandenen
materiellen Dingen mitteilen können.
\par\noindent Beide Bemerkungen hat schon \textsc{Kant} gemacht, aber man begreift nicht, wie
dieser scharfsinnige Philosoph in der ersten einen Beweis für seine Meinung,
dass der Raum nur Form unserer äusseren Anschauung sei, zu finden glauben
konnte, da die zweite so klar das Gegenteil, und dass der Raum unabhängig
von unserer Anschauungsart eine reelle Bedeutung haben muss, beweiset.«

Der zweite Grund findet sich in dem Briefe an \textsc{Wolfgang Bolyai} vom
6. März 1832 (Werke VIII, S. 224), wo sich \textsc{Gauss} über die eben erhaltene
\textit{Appendix} von \textsc{Johann Bolyai}\footnote{\textsc{Johann Bolyai}, \textit{Appendix scientiam spatii absolute veram exhibens}, erschien zuerst 1832; deutsch in \textsc{P. Stäckel}, \textit{Die beiden Bolyai} II, 1913, S. 185 ff.} ausspricht. \textsc{Johann Bolyai} bezeichnet das
System der vom Parallelenaxiom unabhängigen, absoluten Raumlehre, das
noch eine unbestimmte positive Konstante \(i\) enthält, mit \(S\) und das System
der \textsc{Euklid}ischen Geometrie, in dem jene Konstante \(i=\infty\) wird, mit \(\Sigma\), und
sagt am Schlusse seiner Schrift (a. a. O. S. 216): »Es bliebe endlich übrig,
den Beweis dafür zu geben, dass es unmöglich ist, ohne irgend eine Annahme
zu entscheiden, ob \(\Sigma\) oder irgend ein — und welches \(S\) stattfindet; dies mag
jedoch einer geeigneteren Gelegenheit vorbehalten bleiben«. In bezug auf
diesen Punkt schreibt nun \textsc{Gauss} an \textsc{Wolfgang Bolyai}: »Gerade in der Unmöglichkeit,
zwischen \(\Sigma\) und \(S\) a priori zu entscheiden, liegt der klarste Beweis,
dass \textsc{Kant} Unrecht hatte zu behaupten, der Raum sei nur Form unserer
Anschauung. Einen andern, ebenso starken Grund habe ich in einem kleinen
Aufsatze angedeutet, der in den Göttingischen gelehrten Anzeigen 1831 steht«.

Beiden Äusserungen gemein ist die Bezugnahme auf \textsc{Kant}; während jedoch
in der zweiten, an \textsc{Bolyai} gerichteten, diese Bezugnahme rein polemisch
ist\footnote{Auf diese, der Natur der Sache nach wesentlich dialektische Polemik gehen wir hier nicht ein, sondern verweisen auf die Ausführungen bei \textsc{P. Mansion}, \textit{Gauss contre Kant sur la Géométrie Non Euclidienne}, Verhandlungen des III. Intern. Kongr. für Philosophie, Heidelberg 1908, S. 438—447; \textsc{H. E. Timerding}, \textit{Kant und Gauss}, Kant-Studien 28, 1923, S. 16—40; \textsc{A. Galle}, \textit{Gauss und Kant}, Das Weltall 24, 1925, S. 194—200, 230. Im Text beschränken wir uns auf die spezifisch mathematischen Erörterungen.}, beruft sich die erste auf \textsc{Kant} als den Urheber der beiden vorgebrachten
Bemerkungen. In der Tat finden sich diese, im wesentlichen übereinstimmend
\par\clearpage\SourceLeaf{264}\ArithmeticHead{179}{R}{VIIIc. KONFORME ABBILDUNG. GEOMETRISCHES.}
\noindent mit \textsc{Gauss}, in der Abhandlung \textit{Von dem ersten Grunde des Unterschieds der
Gegenden im Raume}, 1768\footnote{\textsc{Kant}s sämtliche Werke herausgeg. von \textsc{Rosenkranz} und \textsc{Schubert}, V, 1839, S. 291 ff.} und in der lateinischen Inauguraldissertation von
1770\footnote{Ebenda, I, 1838, S. 301 ff.}, und zwar mit besonderer Betonung der Existenz von symmetrischen
Körpern, wie z. B. rechte und linke Hand, deren Unterschiede sich in Eigenschaften
kundgeben, die für einen Körper »nicht lediglich auf dem Verhältnis
der Lage seiner Teile gegeneinander beruhen, sondern \mbox{.\,.\,.} auf einer Beziehung
gegen den allgemeinen Raum \mbox{.\,.\,.}, doch so, dass [diese Beziehung] nicht unmittelbar
kann wahrgenommen werden, aber wohl diejenigen Unterschiede
der Körper, die einzig und allein auf diesem Grunde beruhen« (\textsc{Kant}s Werke
V, S. 298). Diese Eigenschaften können deshalb »nur durch Gegenhaltung
mit andern Körpern« aufgezeigt (ebenda, S. 301), nicht aber »diskursiv beschrieben,
d. h. auf verstandesmässige Begriffe zurückgeführt werden« (\textsc{Kant}s
Werke I, S. 321, 322). — Die von \textsc{Gauss} hervorgehobene erste Tatsache
(oben 1. Bemerkung), dass nach willkürlicher Festsetzung von vorwärts—rückwärts
und oben—unten, der Gegensatz rechts—links völlig bestimmt sei, erscheint
ihm offenbar als durch die Erfahrung gegeben. Bei \textsc{Kant} wird diese
Tatsache (\textsc{Kant}s Werke I, S. 321) durch die damit völlig äquivalente Aussage
ersetzt, dass der Raum drei Dimensionen habe; die dreifache Abmessung
des Raumes, von der \textsc{Kant}, z. B. im § 9 der Abhandlung \textit{Gedanken von der
wahren Schätzung der lebendigen Kräfte etc.} (1747, \textsc{Kant}s Werke V, S. 25) sagt,
dass der Grund hierfür noch unbekannt sei, war also für \textsc{Gauss} eine Erfahrungstatsache,
was auch noch dadurch bestätigt werden kann, dass es ja nur durch die
die Erfahrung gegeben ist, ob wir die Lage eines Punktes im Raum durch drei
Abmessungen oder durch mehr oder weniger zu bestimmen imstande sind. Dagegen
wird, wenn dies erst einmal feststeht, die zweite Bemerkung ein beweisbarer
Satz. Diese Bemerkung ist nämlich mit der gleichwertig, die \textsc{Gauss} in dem
Briefe an \textsc{Gerling} vom 23. Juni 1846 (Werke VIII, S. 248) dahin formuliert,
dass drei von einem Punkte ausgehende, nicht in einer Ebene liegende Gerade
(ein sogenanntes Dreibein) mit drei andern gleiche oder verkehrte Lage haben
können, und dass man »diesen Unterschied \mbox{.\,.\,.} nicht auf Begriffe bringen,
sondern nur \mbox{.\,.\,.} an wirklich vorhandenen räumlichen Dingen vorzeigen kann«\footnote{Vergl. diese Äusserung mit den oben zitierten Worten \textsc{Kant}s.}.
\HistoricalFoot{}{23\textsuperscript{*}}
\par\clearpage\SourceLeaf{265}\ArithmeticHead{180}{L}{L. SCHLESINGER, ÜBER GAUSS' FUNKTIONENTHEORETISCHE ARBEITEN.}
\noindent Der Beweis kann durch ähnliche Überlegungen erbracht werden, wie die oben
skizzierten in den Artikeln 8 und 14 der Preisarbeit über die ähnliche und
verkehrte Lage bei der konformen Abbildung; man vergleiche etwa die Entwicklungen
bei \textsc{C. Lange}\begingroup\renewcommand{\thefootnote}{270a)}\footnote{\textsc{C. Lange}, \textit{Beitrag zur analytischen Geometrie der geraden Linie im Raum}, Programm des Gymnasiums \mbox{.\,.\,.\,.} in Insterburg 1864, zitiert bei \textsc{E. Müller}, Encyklopädie, Band III, 1, S. 620.}\addtocounter{footnote}{-1}\endgroup, die in sehr klarer Weise zeigen, dass sich zwar
durch messbare Grössen feststellen lässt, ob zwei Dreibeine gleich oder entgegengesetzt
orientiert sind, nicht aber, welche Orientierung für ein einzelnes
Dreibein stattfindet, indem z. B. alle Formeln der analytischen Geometrie ungeändert
bleiben, wenn man die Orientierung des Koordinatendreibeins umkehrt.
Ohne Beziehung auf ein festes Raumsystem ist es also nicht möglich,
zwei symmetrische Körper in dem Sinne von einander zu unterscheiden\footnote{Vergl. die zitierten Aussprüche von \textsc{Kant}.},
dass man Kennzeichen angibt, die dem einen zukommen, dem andern nicht;
denn in allem, was messbar ist, stimmen zwei symmetrische Körper überein,
und das, worin sie sich unterscheiden, entzieht sich jeder Grössenbestimmung
und damit jeder mathematischen Ausdrucksweise.

Was nun den zweiten der beiden von \textsc{Gauss} vorgebrachten Gründe für
den empirischen Charakter des Raumes anlangt, der in der Briefstelle an
\textsc{Bolyai} ausgesprochen wird, so scheint \textsc{Gauss} die daselbst hervorgehobene
»Unmöglichkeit, a priori zwischen \(\Sigma\) und \(S\) zu unterscheiden«, als erwiesen
angesehen zu haben, d. h. die Widerspruchsfreiheit des Systems \(S\) der absoluten
Geometrie stand, wie bereits (S. 176) bemerkt wurde, für ihn fest. So
schreibt er auch am 8. November 1824 an \textsc{Taurinus} (Werke VIII, S. 187): »Alle
meine Bemühungen einen Widerspruch, eine Inkonsequenz in dieser nichteuklidischen
Geometrie zu finden, sind fruchtlos gewesen«. Wahrscheinlich
hat er aber diese Widerspruchslosigkeit auch noch in ähnlicher Weise begründet,
wie später \textsc{Riemann}, indem er nämlich nicht nur die Ebene, sondern
auch den dreidimensionalen Raum des Systems \(S\) als eine Zahlenmannigfaltigkeit
mit bestimmter Metrik auffasste. Neben diese abstrakte Frage der Unmöglichkeit,
a priori zwischen \(\Sigma\) und \(S\) zu entscheiden, tritt für \textsc{Gauss} aber
noch die, festzustellen, welches die »wahre«, d. h. also die den physikalisch-astronomischen
Messungen zugrunde liegende Geometrie ist. Hierüber heisst
es in dem erwähnten Briefe an \textsc{Taurinus}: »Wäre die nichteuklidische
\par\clearpage\SourceLeaf{266}\ArithmeticHead{181}{R}{VIIId. DIE HUNDERT THEOREME.}
\noindent Geometrie die wahre und jene Konstante [das \textsc{Bolyai}sche \(i\)] in einigem Verhältnis
zu solchen Grössen, die im Bereich unserer Messungen auf der Erde oder am
Himmel liegen, so liesse sie sich a posteriori ausmitteln«; für \textsc{Gauss} war also
die Bestimmung der Metrik des »wahren«, d. h. physikalisch-astronomischen
Raumes eine empirisch zu lösende Aufgabe.

Wir können hiernach die »Metaphysik des Räumlichen«, wie sie sich für
\textsc{Gauss} auf Grund seiner geometrisch-analytischen Einsichten darstellte, wie
folgt charakterisieren: Unter den \(n\)-fach ausgedehnten Zahlenmannigfaltigkeiten
— oder, wie es in der Jubiläumsschrift von 1849 heisst, »nach der Stetigkeit
zusammenhängenden Grössenkombinationen« — ist auch der »wahre Raum«
begriffen. Auf Grund der Erfahrung bestimmt sich einerseits der Wert von
\(n\), d. h. die Dimensionszahl, und andererseits die Metrik, d. h. der angenäherte
Wert der Raumkonstanten, des \textsc{Bolyai}schen \(i\).

»Sollte ich in diesem Leben« — so schreibt \textsc{Gauss} am 28. Juli 1823 mit
Bezug auf die Preisschrift zur konformen Abbildung an \textsc{Olbers} — »noch einmal
in eine dem Arbeiten günstigere Lage kommen, so werde ich diese Abhandlung
mit als Teil einer viel ausgedehnteren Untersuchung verarbeiten«\footnote{\textit{Briefwechsel Olbers-Gauss} I, S. 252.}.
Wir haben im vorstehenden versucht, einen Überblick über den Umfang dieser
an das Problem der konformen Abbildung anschliessenden Untersuchungen zu
gewinnen, die ja auch aufs engste mit \textsc{Gauss}' Auffassung der komplexen
Grössen zusammenhängen. — Die etwas bittere Stimmung, die in den Worten
der angeführten Briefstelle durchklingt, hängt einerseits mit der Überlastung
durch die praktischen Arbeiten im Gelände und am Rechentische zusammen,
die mit der \textsc{Gauss} 1820 übertragenen Gradmessung einhergingen, andererseits
aber wohl auch mit »seinen kollegialischen Verhältnissen, mit dem, der ihm
so nahe ist«, und »gerechten Sorgen für die Zukunft, über die \textsc{Gauss}' Schwiegermutter,
Frau \textsc{Waldeck}, in einem Briefe an \textsc{Olbers} vom 14. März 1821 in beweglichen
Worten klagt\footnote{\textit{Briefe zwischen A. v. Humboldt und Gauss}, 1877, S. 6, 7.}.
